\subsection{7月}

\subsubsection{7月5日 周日 —— 深夜急行}

\begin{itemize}
  \item \textbf{突然接到出差通知}：周日晚上突然收到消息，要求周一赴北京出差。事出紧急，连夜收拾行李乘坐高铁出发。
  \item \textbf{深夜抵京}：到达北京已是深夜，折腾到凌晨2点多才睡着，第二天还要直接进入工作状态，确实有些疲惫。
\end{itemize}

\begin{personalnote}
这次出差来得太突然，从收到消息到坐上高铁不过几个小时。连夜赶路的经历虽然辛苦，但也算是研究生阶段一次难得的"说走就走"的体验。第一次切身感受到项目申报的时间紧迫感。
\end{personalnote}

\subsubsection{7月6日 $\sim$ 7月9日 周一$\sim$周四 —— 北京出差：PLL 申报书撰写}

\begin{itemize}
  \item \textbf{PLL 项目申报书撰写}：在北京的四天全身心投入到 PLL（锁相环）相关的项目申报书撰写工作中。从技术方案论证到申报书框架搭建，再到具体内容填充，密集推进。
  \item \textbf{高强度协作}：与北京方面的团队对接，讨论技术路线和申报策略，期间经历了多轮修改和打磨。
  \item \textbf{周四首飞体验}：周四（7月9日）返程时第一次坐飞机——从北京飞回。虽然是出差结束的正常行程，但也算是解锁了人生第一次乘机体验。
\end{itemize}

\begin{personalnote}
这次北京出差虽然来得突然，但整体收获不小。PLL 申报书的撰写让我对锁相环技术的应用场景和前沿方向有了更系统的认识。周四第一次坐飞机的体验也算是这趟匆忙出差的一个小小纪念——以前总觉得坐飞机是件大事，实际体验下来也就是交通工具的一种罢了（虽然起飞时的推背感还是有点意思）。密集的几天下来，深深体会到项目申报的不易，每一个技术指标都要反复推敲。
\end{personalnote}
